\section{Implementation in Stan}\label{app:stan-implementation}

For monotonicity and convexity, the Stan programs implement
\cref{alg:posterior} using the No-U-Turn sampler (NUTS)
\citep{hoffman2014nuts,carpenter2017stan}. The sampled parameter is
the unconstrained vector $\boldsymbol z\in\mathbb R^m$. Gaussian
preparation supplies $\boldsymbol m_g,\boldsymbol B,m_\gamma,
\boldsymbol\lambda$, and $s>0$ as data, together with the basis
geometry, constraint interval, and threshold tolerance. The threshold
routine also has a specified precision budget. The rank and padding satisfy the conditions of
\cref{cor:monotone,cor:convex}. Preparation allows a general
positive-definite noise covariance. The convexity program uses a basis
centered at zero, so a model centered at $c$ is represented by
subtracting $c$ from both observation and constraint coordinates.

At each target evaluation, the program forms
$\boldsymbol g=\boldsymbol m_g+\boldsymbol B\boldsymbol z$ and
$\mu=m_\gamma+\boldsymbol\lambda^T\boldsymbol z$, then evaluates
the upper threshold $\widetilde\tau_r(\boldsymbol g)$.
In \cref{lst:stan-core}, \texttt{threshold\_upper} denotes this
value, \texttt{conditional\_sd} denotes $s$, and
\texttt{alpha} denotes
$\widetilde\alpha=(\widetilde\tau_r(\boldsymbol g)-\mu)/s$.
The Gaussian inputs already incorporate the observations, so the
model block adds only the log survival weight
\texttt{log\_survival}$=\log\Phi^{\mathrm c}(\widetilde\alpha)$ to the
standard normal log density of $\boldsymbol z$.

External C++ functions supply the threshold and normal-tail operations.
The threshold function follows
\cref{subsec:threshold-computation} and attaches the approximate
envelope derivative from \cref{subsec:threshold-derivatives} to
its returned upper value through a custom automatic differentiation
adapter. Stan can therefore propagate the force through the linear
coefficient map without differentiating the root search or interval
subdivision. The log-survival adapter uses the negative inverse
Mills ratio, and the quantile routine accepts a log probability,
as described in \cref{subsec:tail-computation}.

\begin{listing}[htbp]
\begin{minted}[fontsize=\footnotesize,breaklines]{stan}
model {
  target += -0.5 * dot_self(z) + log_survival;
}
generated quantities {
  real eta = -normal_log_quantile(
    log_survival + log(uniform_rng(0, 1)));
  real gamma =
    threshold_upper + conditional_sd * (eta - alpha);
}
\end{minted}
\caption{Marginal sampling and conditional reconstruction in Stan.
The excerpt assumes that the Gaussian transformations, upper threshold,
and log survival weight have been evaluated. The generic routine
\texttt{normal\_log\_quantile} returns $\Phi^{-1}(\exp(\ell))$
from a supplied log probability $\ell$.}
\label{lst:stan-core}
\end{listing}

For each retained $\boldsymbol z$, the generated-quantities block
draws $\gamma$ from its truncated conditional Gaussian law using an
independent uniform variate and the same upper threshold as the model
block. Inverting the coefficient map then gives a joint coefficient
draw. In particular, the monotonicity program restores $a_2=-\gamma$
and $a_4=g_4-\gamma/4$, while the convexity program restores
$a_1=-\gamma$. The other coefficients follow
\cref{tab:coefficient-reconstruction}. The finite-value and quantile
checks in \cref{subsec:tail-computation} apply to this reconstruction.

Extending these programs to a joint restriction requires Gaussian
inputs and reconstruction under the common map in
\cref{cor:combinations}. The external threshold routine must enclose
$\tau_A^\varepsilon$ and supply the envelope force from its active
orders and locations, as described in
\cref{subsec:threshold-computation,subsec:threshold-derivatives}.
With these inputs, the scalar sampling and reconstruction blocks
in \cref{lst:stan-core} apply unchanged.

The returned log values determine the density used to weight NUTS
states, whereas the supplied force determines the proposed trajectory.
To justify this distinction, consider exact arithmetic with a fixed
Euclidean metric, step size, and maximum tree depth after warm-up.
Suppose the log values and force are finite, deterministic, measurable
functions of position and the resulting density is normalizable.
Leapfrog consists of position and momentum shears that preserve volume
and invert when the step size changes sign. Consequently, momentum
reversal makes the proposal reversible even when the force is not the
log-density gradient \citep[Section 5.5]{neal2011hmc}.
With the recursive subtree checks and weighted state selection of
NUTS \citep[Appendices A.3--A.5]{betancourt2017conceptual}, this
preserves the density defined by the returned log values in the ideal
construction without evaluation failures or divergence truncation.
The depth limit makes each transition finite. This argument justifies
using an approximate envelope force but gives no bound on finite-chain
or floating-point error, and the force approximation can still affect
sampling efficiency.
